% Receive down the left, transmit up the right, one peripheral between them.
\begin{tikzpicture}[font=\sffamily\small]
% ---------------- receive
\node[fw, text width=28mm] (span) at (2.5,3.0) {span from chapter 4};
\node[fw, text width=28mm] (feed) at (2.5,1.5) {frame\_feed()\\{\scriptsize two states}};
\node[fw, text width=28mm] (unst) at (2.5,0.0) {unstuff\\{\scriptsize vendored, MIT}};
\node[fw, text width=28mm] (cmp) at (2.5,-1.5) {crc16()\\{\scriptsize compare}};
\node[user, text width=28mm] (app) at (2.5,-3.0) {application\\{\scriptsize whole messages only}};

% ---------------- transmit
\node[fw, text width=28mm] (enc) at (9.6,-3.0) {frame encode};
\node[fw, text width=28mm] (add) at (9.6,-1.5) {crc16()\\{\scriptsize append}};
\node[fw, text width=28mm] (stf) at (9.6,0.0) {stuff\\{\scriptsize vendored, MIT}};
\node[fw, text width=28mm] (tx) at (9.6,1.5) {USART3 transmit};
\node[ext, text width=28mm] (host) at (9.6,3.0) {host, twin.py};

% ---------------- the only silicon this chapter adds
\node[hw, text width=26mm] (crc) at (6.05,-1.5) {checksum unit\\{\scriptsize POL, INIT, CR}};

\draw[flow] (span) -- (feed);
\draw[flow] (feed) -- (unst);
\draw[flow] (unst) -- (cmp);
\draw[flow] (cmp) -- (app);
\draw[flow] (app) -- (enc);
\draw[flow] (enc) -- (add);
\draw[flow] (add) -- (stf);
\draw[flow] (stf) -- (tx);
\draw[flow] (tx) -- (host);
\draw[bus] (cmp) -- (crc);
\draw[bus] (crc) -- (add);

\node[umlnote, text width=58mm, anchor=north west] at (0.9,-4.4)
  {One checksum routine serves both directions, behind one name chosen at build
   time. Two implementations in one program means only one of them is tested,
   and it is never the one that fails.};
\end{tikzpicture}
